\documentclass[10pt,a4paper]{amsart}
\usepackage[margin=19mm]{geometry}
\usepackage{amsmath,amssymb,mathtools}
\usepackage[scr=rsfs,cal=euler]{mathalfa}
\usepackage{xfrac}
\usepackage[unicode,hidelinks,bookmarks=false]{hyperref}
\newcommand{\ie}{i.e.\ }
\newcommand{\cf}{cf.\ }
\newcommand{\resp}{respectively\ }
\newcommand{\Subsection}{\subsection}
\providecommand{\nobreakdash}{\nobreak\hbox{-}}
\setlength{\emergencystretch}{2em}
\multlinegap0pt
\usepackage{mathtools}
\usepackage{amssymb}
\usepackage{dsfont}
\usepackage{xcolor}
\usepackage{tikz}
\newtheorem{theorem}{Theorem}
\newtheorem{lemma}{Lemma}
\usepackage{cleveref}
\crefname{theorem}{Theorem}{Theorems}
\crefname{lemma}{Lemma}{Lemmas}
\newcommand{\den}[2]{p(#1;\,#2)}
\newcommand{\poshom}[1]{\mathrm{Hom}^+(\mathcal{A}^{#1}, \mathbb{R})}
\newcommand{\tdensity}[2]{\pi(#1;\,#2)}
\title{Formalizing Flag Algebras in Lean}
\author{Gyeongwon Jeong, Seonghun Park, Jihoon Hyun, Sang-il Oum, Hongseok Yang}
\date{Source: arXiv:2607.23500v1 (CC BY 4.0)}
\begin{document}
\maketitle

\noindent\emph{Source and licence.} Formalizing Flag Algebras in Lean. Version arXiv:2607.23500v1 (CC BY 4.0), distributed by arXiv under the Creative Commons Attribution 4.0 International licence, http://creativecommons.org/licenses/by/4.0/, as recorded in the arXiv licence metadata for this exact version and shown on the abstract page https://arxiv.org/abs/2607.23500v1. Full licence notice: \texttt{licenses/arxiv-2607.23500-CC-BY-4.0.txt}.\par\smallskip

\noindent\textit{Editorial scope.} Counted unit: the article's lemma lem:semantic-bound (source0.tex lines 931--934). The article's complete original proof of it is delivered with it (source0.tex lines 935--948, one \texttt{proof} environment of 14 lines). No mathematics is written, repaired or re-derived: the statement and the proof are the article's own lines, in the article's own notation, and no retained line is substituted or re-flowed.  Retained uncounted as the source-local context the proof needs: lines 881--892.  Nothing was omitted from any retained span, so the package carries no omission record.  No retained line contains a citation, so the wrapper carries no bibliography and no bibliography entry.  No retained line contains a figure environment, an image-inclusion command or drawing code, so no image is delivered and the package declares no asset dependency.  Editorial formatting only, in addition: an AMS \texttt{amsart} preamble that restates the article's own declarations and macros used by the retained lines, namely the environments \texttt{theorem}, \texttt{lemma} on separate counters, together with the standard packages those lines need.  The retained environments print in this document as Theorem 1, Lemma 1 in order of appearance, whereas the article prints the same items with the numbering of its own full document.  The complete native source of the article as distributed in the arXiv e-print for this exact version is bundled unchanged as \texttt{sources/206032/source0.tex}.
\begin{theorem}[Razborov]\label{thm:convergent-hom}
  \begin{enumerate}
  \item[\textup{(a)}] If $\{G_n\}_{n\in\mathbb{N}}$ is a convergent sequence of
    $\sigma$-flags, then
    \[
      \lim_{n\to\infty} p^{G_n} \in \poshom{\sigma}.
    \]
  \item[\textup{(b)}] Conversely, every $\phi \in \poshom{\sigma}$ satisfies
    $\phi = \lim_{n\to\infty} p^{G_n}$ for some convergent sequence
    $\{G_n\}_{n\in\mathbb{N}}$ of $\sigma$-flags.
  \end{enumerate}
\end{theorem}

\begin{lemma}\label{lem:semantic-bound}
  If $[F] \leq_{\mathcal{H}} c \cdot \mathbf{1}$ for some $c \in \mathbb{R}$, then
  $\tdensity{F}{\mathcal{H}} \leq c$.
\end{lemma}

\begin{proof}
  Suppose for contradiction that $\tdensity{F}{\mathcal{H}} > c$.  Then
  there exist $\varepsilon > 0$ and an increasing sequence of
  \(\mathcal{H}\)-free graphs $\{G_k\}_{k\in\mathbb{N}}$ such that
  $\den{F}{G_k} > c + \varepsilon$ for all $k$.  Since each $p^{G_k}$
  is a point in $[0,1]^{\mathcal{F}^\emptyset}$ and
  $\mathcal{F}^\emptyset$ is countable, this product space is compact and
  metrizable.  Hence the sequence $\{p^{G_k}\}$ has a convergent
  subsequence $\{p^{G_{k_j}}\}$; by \Cref{thm:convergent-hom}(a) its limit is some
  $\phi \in \poshom{\emptyset}$.  Since $\phi$ is the pointwise limit of
  $\{p^{G_{k_j}}\}$, we have
  $\phi([F]) = \lim_j \den{F}{G_{k_j}} \geq c + \varepsilon > c$.  But
  $[F] \leq_{\mathcal{H}} c \cdot \mathbf{1}$ means $\phi([F]) \leq c$ for every~$\phi \in \poshom{\emptyset}$, a contradiction.
\end{proof}

\end{document}
